\subsection{截面}

7.4.1. 设 M 是以 B 为基底的向量丛。对 B 的每个开集 U，记 \(\mathcal{S}_{\mathbf{M}}^{r}(\mathbf{U})\) 为 M 在 U 上的类为 \(C^r\) 的截面集，即从 U 到 M 的态射 \(s\)，其类为 \(C^r\)，满足 \(s(b) \in M_b\) 对所有 \(b \in \mathbf{U}\) 都成立。根据以下规则，这个集合带有关于态射函数环 \(\mathcal{C}^r(\mathbf{U})\) 的模结构：

\begin{enumerate}
\def\labelenumi{(\arabic{enumi})}
\tightlist
\item
\end{enumerate}

\[
(s + s ^ {\prime}) (b) = s (b) + s ^ {\prime} (b)\tag{2}
\]

\[
(\varphi . s) (b) = \varphi (b). s (b)
\]

其中 \(s, s'\) 属于 \(\mathcal{S}_{\mathrm{M}}^{r}(\mathrm{U})\)，\(\varphi\) 属于 \(\mathcal{C}^{r}(\mathrm{U})\)。当开集 U 变化时，得到一个从 B 到 M 的映射层 \(\mathcal{S}_{\mathrm{M}}^{r}\)（参见第 5.4.1 号），称为 M 的截面层。

7.4.2. 设 \(M_{1},\ldots,M_{d}\) 和 N 是以 B 为基底的向量丛，u 是从 \(M_{1}\times_{B}\ldots\times_{B}M_{d}\) 到 N 的多线性态射。对 \(1\leqslant j\leqslant d\)，给定一个截面 \(s_{j}\)，它是 \(M_{j}\) 在 B 的一个开集 U 上的截面；按下式在 U 上定义 N 的截面 \(u(s_{1},\ldots,s_{d})\)：

\[
u (s _ {1}, \dots , s _ {d}) (b) = u _ {b} (s _ {1} (b), \dots , s _ {d} (b)) \quad \text {   for   } b \in \mathrm{U}.
\]

映射 \((s_1, \ldots, s_d) \mapsto u(s_1, \ldots, s_d)\) 是 \(\mathcal{C}^r(\mathrm{U})\)-多线性的，有时记为 \(\mathcal{S}(u)\)。

7.4.3. 设 \(f\) 是从流形 \(\mathbf{B}'\) 到 \(\mathbf{B}\) 的态射，\(\mathbf{M}\) 是以 \(\mathbf{B}\) 为基底的向量丛。对每个开集 \(\mathbf{U}\)，它属于 \(\mathbf{B}\)，以及每个 \(s \in \mathcal{S}_{\mathbf{M}}^{r}(\mathbf{U})\)，映射 \(x \mapsto (x, s(f(x)))\) 是类为 \(\mathbf{C}^r\) 的 \(f^{*}\mathbf{M}\) 在开子集 \(f^{-1}(\mathbf{U})\) 上的截面，记为 \(f^{*}s\)，称为 \(s\) 关于 \(f\) 的逆像。映射 \(s \mapsto f^{*}s\) 从 \(\mathcal{S}_{\mathbf{M}}^{r}(\mathbf{U})\) 到 \(\mathcal{S}_{f^{*}\mathbf{M}}^{r}(f^{-1}(\mathbf{U}))\) 关于同态 \(g \mapsto g \circ (f|f^{-1}(\mathbf{U}))\)（它从 \(\mathcal{C}^{r}(\mathbf{U})\) 到 \(\mathcal{C}^{r}(f^{-1}(\mathbf{U}))\)）是半线性的。

此外，若 N 是以 B' 为基底的向量丛，g 是从 M 到 N 的 f-余态射，有时将记为 \(\mathcal{S}(g)\) 的映射 \(s \mapsto g \circ f^{*}s\) 从 \(\mathcal{S}_{\mathbf{M}}^{\mathbf{r}}(\mathbf{U})\) 到 \(\mathcal{S}_{\mathbf{N}}^{\mathbf{r}}(f^{-1}(\mathbf{U}))\)。

7.4.4. 设 M 是以 B 为基底的有限秩向量丛。M 在 B 的一个开集 U 上的标架定义为 M 在 U 上的截面组成的有限序列 \((s_{1},\ldots,s_{n})\)，使得 \((s_{1}(b),\ldots,s_{n}(b))\) 是向量空间 \(M_{b}\) 的一个基，对所有 \(b\in U\) 均如此。于是序列 \((s_{1},\ldots,s_{n})\) 是 \(\mathcal{C}^{r}(\mathrm{U})\)-模 \(\mathcal{S}_{\mathbf{M}}^{r}(\mathrm{U})\) 的一个基。若 f 是从流形 \(B'\) 到 B 的态射，则截面 \(f^{*}s_{j}\) 构成 \(f^{*}\mathbf{M}\) 在 \(f^{-1}(\mathbf{U})\) 上的一个标架。

7.4.5. 设 L 是一个赋有有限维 K-代数结构的域，(M, B, π) 是一个纤维化。假定每个纤维 \(M_{b}\) 上都给定了有限维的 L-向量空间结构。M 上至多存在一个以 B 为基底的向量丛结构，它与映射 π、M 的流形结构以及各纤维上的 L-向量空间结构相容（7.3.4）。这种结构存在的充要条件是满足以下条件：

（FV）对每个 \(b_0 \in \mathbf{B}\)，存在一个开邻域 \(\mathbf{U}\)，其中包含 \(b_0\) 且位于 \(\mathbf{B}\) 中，以及一个流形同构 \(\varphi\)，它从 \(\pi^{-1}(\mathbf{U})\) 到乘积 \(\mathbf{U} \times \mathbf{F}\)，它是 \(\mathbf{U}\) 与 \(\mathbf{F}\) 的乘积，其中后者是有限维的 \(\mathbf{L}\)-向量空间，使得对每个 \(b \in \mathbf{U}\)，双射 \(\varphi_b\) 从 \(\mathbf{M}_b\) 到 \(\mathbf{F}\) 由 \(\varphi\) 诱导，并且是 \(\mathbf{L}\) 上的向量空间同构。

于是，满足（FV）的三元组（U，\(\varphi\)，F）就是向量丛 M 的向量丛图。

条件（FV）等价于：

（FV'）对每个 \(b_0 \in \mathbf{B}\)，存在一个整数 \(n\) 以及 \(n\) 个截面 \(s_1, \ldots, s_n\)，它们是 \(\mathbf{M}\) 在开邻域 \(\mathbf{U}\) 上的截面，而该邻域包含 \(b_0\)，使得映射

\[
(b, a _ {1}, \dots , a _ {n}) \mapsto a _ {1} s _ {1} (b) + \dots + a _ {n} s _ {n} (b)
\]

是从流形 \(\mathbf{U} \times \mathbf{L}^{n}\) 到流形 \(\pi^{-1}(\mathbf{U})\) 的同构。

7.4.6. 设 \(M_{1},\ldots,M_{d}\) 和 N 是以 B 为基底的向量丛，其中 \(M_{j}\) 为有限秩。假定对 B 的每个开集 U，都给定了一个 \(\mathcal{C}^{r}(\mathrm{U})\)-多线性映射 \(\varphi_{\mathrm{U}}\)，从 \(\mathcal{S}_{\mathrm{M}_{1}}^{\mathrm{r}}(\mathrm{U})\times\cdots\times\mathcal{S}_{\mathrm{M}_{d}}^{\mathrm{r}}(\mathrm{U})\) 到 \(\mathcal{S}_{\mathrm{N}}^{\mathrm{r}}(\mathrm{U})\)，使得当 \(V \subset U\) 时有：

\[
\varphi_ {\mathrm{U}} (s _ {1}, \dots , s _ {d}) | \mathrm{V} = \varphi_ {\mathrm{V}} (s _ {1} | \mathrm{V}, \dots , s _ {d} | \mathrm{V}).
\]

于是存在一个多线性态射 \(u\)，从 \(\mathbf{M}_1 \times_{\mathbf{B}} \ldots \times_{\mathbf{B}} \mathbf{M}_d\) 到 \(\mathbf{N}\)，使得 \(\varphi_{\mathbf{U}}(s_1, \ldots, s_d) = u(s_1, \ldots, s_d)\) 对任意截面 \(s_j\)，属于 \(\mathbf{M}_j\) 在开集 \(\mathbf{U}\)（属于 \(\mathbf{B}\)）上的截面，均成立。

7.4.7. 设 \(f\) 是从流形 \(\mathbf{B}'\) 到 \(\mathbf{B}\) 的态射，且 \(\mathbf{M}\)（分别 \(\mathbf{M}'\)）是以 \(\mathbf{B}\)（分别 \(\mathbf{B}'\)）为基底的有限秩向量丛。假定对每个开集 \(\mathbf{U}\)，它属于 \(\mathbf{B}\)，给定一个连续的 \(\mathcal{C}^{r}(\mathbf{U})\)-半线性映射 \(\varphi_{\mathbf{U}}\)，从 \(\mathcal{S}_{\mathbf{M}}^{r}(\mathbf{U})\) 到 \(\mathcal{S}_{\mathbf{M}'}^{r}(f^{-1}(\mathbf{U}))\)，使得

\[
\varphi_ {\mathbf {U}} (s) | f ^ {- 1} (\mathbf {V}) = \varphi_ {\mathbf {V}} (s | \mathbf {V})
\]

对每个开集 \(V \subset U\) 都有。于是存在唯一一个从 M 到 \(M'\) 的 f-余态射 g，使得 \(\varphi_{\mathrm{U}}(s) = \mathcal{S}(g)(s)\) 对所有 \(s \in \mathcal{S}_{\mathrm{M}}^{r}(\mathrm{U})\) 都成立。

*7.4.8. 设 \(\mathcal{F}\) 是关于环层 \(\mathcal{C}_{\mathrm{B}}^{r}\) 的模层。我们称 \(\mathcal{F}\) 是局部自由的：对每个 \(b \in \mathrm{B}\)，存在一个开邻域 U，含有 \(b\)，以及一个整数 \(n\)，使得 \(\mathcal{F}|\mathrm{U}\) 与一个 \(\mathcal{C}_{\mathrm{U}}^{r}\)-模层 \((\mathcal{C}_{\mathrm{U}}^{r})^{n}\) 同构。

若 M 是以 B 为基底的有限秩向量丛，则层 \(\mathcal{S}_{\mathrm{M}}^{r}\) 是局部自由的。反之，对 B 上每个局部自由层 \(\mathcal{F}\)，存在一个向量丛 M 以及从 \(\mathcal{S}_{\mathrm{M}}^{r}\) 到 \(\mathcal{F}\) 的层同构。若 M' 是另一个以 B 为基底的有限秩向量丛，则映射 \(g \mapsto \mathcal{S}(g)\) 是从 M 到 M' 的 B-态射集到从 \(\mathcal{C}^{r}\)-模层态射 \(\mathcal{S}_{\mathrm{M}}^{r}\) 到 \(\mathcal{S}_{\mathrm{M}'}^{r}\) 的集合的双射。

